\begin{table}
\begin{tcolorbox}[tabularx={X|X|},enhanced,fonttitle=\bfseries\large,fontupper=\normalsize\sffamily,
colback=yellow!10!white,colframe=green!50!black,colbacktitle=blue!30!white,
coltitle=black,title=Mathematical Data in Orbiter]
\hline
Object & Description \\
\hline
\hline
BLT-sets & BLT sets of $Q(4,q)$ exist for all odd prime powers. 
The classification of BLT-sets of $Q(4,q)$ is known to Orbiter for all $q \le 73.$\\
\hline
Cubic Surfaces & Cubic surfaces with 27 lines exist 
for all finite fields apart from ${\mathbb F}_2,$ ${\mathbb F}_3,$ ${\mathbb F}_5.$ 
Orbiter knows the classification of cubic surfaces with 27 lines 
for all fields ${\mathbb F}_q$ of order $q \le 128.$  \\
\hline
Quartic curves & Orbiter knows the classification of smooth quartic curves with 28 bitangents 
in projective planes over all fields ${\mathbb F}_q$ for $q=9,$ $13,$ $19,$ $25,$ $27,$ $29,$ $31.$ \\
\hline
Spreads & A spread is a set of $q^k+1$ pairwise non-intersecting $k$-dimensional subspaces 
of $\bbF_q^{2k}.$ 
Spreads are related to translation planes of order $q^k$.
Orbiter knows the classification of spreads for 
$(q,k) \in \{(2,2), (3,2), (2,4), (4,2), (5,2), (3,3)\}.$ \\
\hline
Hyperovals & A hyperoval in $\PG(2,2^e)$ is a set of $2^e+2$ points, no three collinear.
Orbiter knows the classification of hyperovals for $e=3,4,5.$\\
\hline
Dual hyperovals & A $k$-dimensional dual hyperoval in an ambient space ${\mathbb F}_2^n$ is called a DH$(k,n).$ Orbiter knows the classification of dual hyperovals DH$(4,7)$ and  DH$(4,8).$ \\
\hline
Packings & Orbiter knows the classification of packings of $\PG(3,3).$\\
\hline
\end{tcolorbox}
\caption{\label{tab:mathdata}Mathematical Data in Orbiter}
\end{table}
